\scanpage{28}
\mspage{27}
On doit donc poser

\uline{A\deleted{\unsure{C}}\unsure{E|t|C}7}\quad Si $X,Y,Z\in M$ sont tels que
\[
 X\overset{\circ}{\ll}Y\overset{\circ}{\ll}Z\qquad\text{alors}\qquad X\overset{\circ}{\ll}Z
\]
i.e. la relation $\overset{\circ}{\ll}$ est \uline{transitive}.

Ceci \unsure{assure} donc $\Rightarrow$ dans la prop.
En sens inverse, S'il \unsure{existe|exist} $f$ comme dans (1.42), soit $X\in\widetilde G$, \deleted{alors $\varphi_{\unsure{G|C}}(\unsure{X|x}$} \deleted{\unsure{on}} on sait que $\varphi_G(x)=X$ (tautologique\deleted{\unsure{?}}), moyennant (6), \deleted{on a alors} soit $f(x)=Y$, on a alors $Y=\varphi_F(x)$ i.e. $Y$ est le plus petit élément de $\widetilde F$ \deleted{\unsure{contenant} $X$}. \deleted{\unsure{Donc}} tel que $X\ll Y$. Cela \deleted{signifie} implique donc que pour tt $X\in\deleted{\unsure{\widetilde F}}\widetilde G$, \unsure{existe|exist} $Y\in\deleted{\unsure{\widetilde G}}\widetilde F$ telle que $X\ll Y$, i.e. $\deleted{\unsure{\widetilde G|\widetilde F}}$ (C4) $\widetilde G\ll\widetilde F$, \unsure{qed.|cqfd.} On n'a pas \unsure{utilisé|utilisé-} en fait que $f$ soit croissante, donc

\uline{Corollaire à la prop.}\quad \deleted{La condition,} \inserted{Pour} que $F\ll G$ i.e. $\unsure{\mathrm{Mulom}|\mathrm{Mulomb}}(F)\ll\mathrm{Mulomb}(G)$ il s. que $\mathrm{Omb}\,F\subset\mathrm{Omb}\,G$,
% slanted
\margin{Ce ne sera sans doute plus vrai, en \unsure{passant} aux \uline{petites} multifigures}
et que \deleted{\unsure{toute classe}} la relation d'équivalence \inserted{$\unsure{\mathcal R_F|R_F}$} sur $\mathrm{Omb}\,F$ définie par $\varphi_F$ soit plus fine que la relation d'équiv. induite \unsure{par|sur} $\mathrm{Omb}\,G$.
